\section{REINFORCE++}



REINFORCE++ modifies PPO in the LLM text by introducing
\begin{itemize}
    % \item using outcome reward model $r(x,y)$, where $x$ and $y$ are input prompt and response, respectively;
    \item token-level KL penalty  $\operatorname{KL}(t;\theta)=\log\frac{\pi_{\theta}(a_t|s_t)}{\pi_{\theta_{\text{ref}}}(a_t|s_t)}$, where $\pi_{\theta_{\text{ref}}}$ can be the SFT model;
    \item token-level reward function $r_t=\Imbb(s_t=[\text{EOS}])r(x,y)-\beta\operatorname{KL}(t)$, where $x$ and $y$ are input prompt and response, respectively and $r(x,y)$ is the outcome reward.
\end{itemize}

Recall the advantage used in PPO is $A^{\operatorname{GAE}(\gamma, \lambda)}_{\phi} = \sum_{\ell=0}^{T-1-t}(\gamma\lambda)^{\ell}\left(r_{t+\ell} + \gamma V^{\pi}_{\phi}(s_{t+1+\ell}) - V^{\pi}_{\phi}(s_{t+\ell}) \right)$. If we drop the value network and plug in the new definition of $r_t$, we get 

\begin{align*}
    A^{\operatorname{GAE}(\gamma, \lambda)}_{t} 
    & = \sum_{\ell=0}^{T-1-t}(\gamma\lambda)^{\ell}\left[\Imbb(s_{t+\ell}=[\text{EOS}])r(x,y)-\beta\operatorname{KL}(t+\ell;\theta) \right] \\
    & = \sum_{i=t}^{T-1}(\gamma\lambda)^{\ell}\left[\Imbb(s_{i}=[\text{EOS}])r(x,y)-\beta\operatorname{KL}(i) \right]  = (\gamma\lambda)^{T-1}r(x,y)-\sum_{i=t}^{T-1}(\gamma\lambda)^{i}\left[\beta\operatorname{KL}(i;\theta) \right].
\end{align*}
Specifically, letting $\gamma=\lambda=1$, we get the advantage estimation used in REINFORCE++ as
\begin{equation}
A^{\operatorname{RF++}}_{t}=r(x,y) - \sum_{i=t}^{T-1}\beta\operatorname{KL}(i;\theta). 
\end{equation}


Consequently, for the $k$th episode, REINFORCE++ solves the following problem

\begin{align*}
     \theta_{k+1}=\arg\max_{\theta} \frac{1}{N}\sum_{n=1}^{N}\frac{1}{T_n}
      \sum_{t=0}^{T_n-1}  
     \min\left\{\frac{\pi_{\theta}(a_t^{(n)}|s_t^{(n)})}{\pi_{\theta_{k}}(a_t^{(n)}|s_t^{(n)})}\bar A_{n,t}, \operatorname{clip}\left(\frac{\pi_{\theta}(a_t|s_t)}{\pi_{\theta_{k}}(a_t|s_t)}, 1-\epsilon, 1+\epsilon\right)
     \bar A_{n,t}\right\},
\end{align*}
where $\bar A_{n,t} = \frac{A^{\operatorname{RF++}}_{n,t}-\mu_k}{\sigma_k}$, $\mu_k=\operatorname{mean}\{A^{\operatorname{RF++}}_{n,t}\}_{n\in[1,N], t\in[0,T_n-1]}$, $\sigma_k=\operatorname{std}\{A^{\operatorname{RF++}}_{n,t}\}_{n\in[1,N], t\in[0,T_n-1]}$, $A^{\operatorname{RF++}}_{n,t}=r(x^{(n)},y^{(n)}) - \sum_{i=t}^{T_n-1}\beta\operatorname{KL}(i;\theta_k)$, and $(x^{(n)},y^{(n)})$ is the $n$th sample.

To conclude, we summarize REINFORCE++ in Algorithm~\ref{algo:RF++}.
\begin{algorithm}[H]
\caption{REINFORCE++}
\begin{algorithmic}[1]
    \State Inputs: Initial Policy network parameter $\theta_0$.
    \For{k=0,1,2....}
        \State Collect a set of trajectories $D_k=\{\tau_n\}_{n=1}^{N}$.
        \State $\{A_{n,t}^{\text{RF++}}\}_{n\in[1,N], t\in[0,T_n-1]}$ calculation, batch mean $\mu_k$ and variance $\sigma_k$  calculation.
        \State Compute normalized advantage estimation $\{\bar A_{n,t} = \frac{A^{\operatorname{RF++}}_{n,t}-\mu_k}{\sigma_k}\}_{n\in[1,N], t\in[0,T_n-1]}$.
        \State Approximately solve 
            \begin{align*}
                 \theta_{k+1}=\arg\max_{\theta} \frac{1}{N}\sum_{n=1}^{N}\frac{1}{T_n}
                  \sum_{t=0}^{T_n-1}  
                 \min\left\{\frac{\pi_{\theta}(a_t^{(n)}|s_t^{(n)})}{\pi_{\theta_{k}}(a_t^{(n)}|s_t^{(n)})}\bar A_{n,t}, \operatorname{clip}\left(\frac{\pi_{\theta}(a_t|s_t)}{\pi_{\theta_{k}}(a_t|s_t)}, 1-\epsilon, 1+\epsilon\right)
                 \bar A_{n,t}\right\},                 
            \end{align*}
    \EndFor
\end{algorithmic}
\label{algo:RF++}
\end{algorithm}


\textbf{Remark}
\begin{itemize}
    \item Compared to GRPO, the normalization happens on the "batch" level; while GRPO applies the normalization within the "group" level.
    \item The KL penalty applied in the loss function is different. For GRPO, the KL penalty is taken out of the definition of the advantage estimation and uses the $k_3$ estimator (recall the definition in Definition\ref{def:kl-est}), which is 
    $\beta\mathbb{D}_{K L}\left[\pi_\theta| | \pi_{\text{ref}}\right](i,t)$
   . In REINFORCE++, the $\operatorname{KL}$ penalty is factorized in the definition of the advantage estimation and uses the $k_1$ estimator, which is  $\beta\sum_{i=t}^{T-1}\operatorname{KL}(i;\theta)$. Notice that in the $t$th action, the penalty is the accumulation of all $\operatorname{KL}$ penalties starting from the time step $t$ to the end.
\end{itemize}

Motivated by the discussion above, \citet{xie2025logic} modifies the loss function by combining the GRPO and REINFROCE++ as below.

 \begin{align*}
     \theta_{k+1}=\arg\max_{\theta} \frac{1}{N}\sum_{n=1}^{N}\frac{1}{T_n}
      \sum_{t=0}^{T_n-1}  
     \min\left\{\frac{\pi_{\theta}(a_t^{(n)}|s_t^{(n)})}{\pi_{\theta_{k}}(a_t^{(n)}|s_t^{(n)})}\bar A_{n,t}, \operatorname{clip}\left(\frac{\pi_{\theta}(a_t|s_t)}{\pi_{\theta_{k}}(a_t|s_t)}, 1-\epsilon, 1+\epsilon\right)
     \bar A_{n,t} -\beta \mathbb{D}_{K L}\left[\pi_\theta| | \pi_{\text{ref}}\right](i,t)\right\},                 
\end{align*}
where $\bar A_{n,t} = \frac{A^{\operatorname{Logic}}_{n,t}-\mu_k}{\sigma_k}$, $\mu_k=\operatorname{mean}\{A^{\operatorname{Logic}}_{n,t}\}_{n\in[1,N], t\in[0,T_n-1]}$, $\sigma_k=\operatorname{std}\{A^{\operatorname{Logic}}_{n,t}\}_{n\in[1,N], t\in[0,T_n-1]}$, $A^{\operatorname{Logic}}_{n,t}=r(x^{(n)},y^{(n)})$. Basically, it moves $\operatorname{KL}$ penalty out of the definition of the advantage estimators and use $k_3$ $\operatorname{KL}$  estimation.


